% Preamble of the GR lecture script. Based on ../template_exercises/template_exercises.tex
% (same packages, logos and header), extended by what a script needs:
% numbered equations per section, exercise/remark boxes, TikZ figures.

% Formatting and layout packages
\usepackage{a4wide}
\usepackage{fancyhdr} % For headers and footers
\usepackage{graphicx} % For including logos and images
\usepackage{parskip} % For better paragraph spacing
\usepackage[utf8]{inputenc}
\usepackage[T1]{fontenc}
\usepackage{amsmath} % Advanced math features
\usepackage{amssymb} % Math symbols
\usepackage{bm} % Bold math symbols
\usepackage{siunitx} % SI units
\usepackage{booktabs} % Better table formatting
\usepackage{subcaption} % Subfigure captions
\usepackage{float} % For float positioning (required for [H])
\usepackage{xcolor}
\usepackage{caption} % For custom caption formatting
\usepackage{geometry} % Better control over margins
\usepackage{enumitem}
\usepackage{tikz}
\usetikzlibrary{arrows.meta, angles, quotes, calc, decorations.pathreplacing}
\usepackage[most]{tcolorbox}
\usepackage[colorlinks=true, linkcolor=blue!50!black, urlcolor=blue!50!black, citecolor=blue!50!black]{hyperref}
\usepackage[nameinlink, capitalise]{cleveref}

\geometry{top=30mm, left=25mm, right=25mm, bottom=30mm, headheight=22mm, headsep=6mm}

% SI unit settings
\sisetup{uncertainty-mode=separate}
\sisetup{exponent-product=\cdot}

% Header and footer (logos as in the exercise template)
\pagestyle{fancy}
\fancyhf{}
\fancyhead[L]{\includegraphics[width=2.5cm]{figures/univie.png}}
\fancyhead[R]{\includegraphics[width=2.2cm]{figures/astro.png}}
\fancyhead[C]{\small\textit{General Relativity -- Lecture Notes}}
\fancyfoot[C]{\thepage}
\renewcommand{\headrulewidth}{0.4pt}

\setlength{\parindent}{0pt}
\numberwithin{equation}{section}
\numberwithin{figure}{section}
\captionsetup{font=small, labelfont=bf}

% ---------- Boxes ----------
% Exercises are numbered within the section so that later chapters can refer to them.
\newtcolorbox[auto counter, number within=section, crefname={Exercise}{Exercises}]{exercise}[1][]{%
  enhanced, breakable, colback=blue!4, colframe=blue!45!black, boxrule=0.5pt, arc=1.5pt,
  left=6pt, right=6pt, top=4pt, bottom=4pt,
  fonttitle=\bfseries, title={Exercise \thetcbcounter}, #1}

\newtcolorbox{remark}[1][Remark]{%
  enhanced, breakable, colback=black!4, colframe=black!45, boxrule=0.4pt, arc=1.5pt,
  left=6pt, right=6pt, top=4pt, bottom=4pt,
  fonttitle=\bfseries, title={#1}}

% Key results: \begin{keyresult} ... \end{keyresult}
\newtcolorbox{keyresult}{%
  enhanced, colback=orange!8, colframe=orange!60!black, boxrule=0.5pt, arc=1.5pt,
  left=6pt, right=6pt, top=2pt, bottom=2pt}

% Marks where a lecture of the course starts: \lecture{3}
\newcommand{\lecture}[1]{%
  \par\medskip\noindent
  {\color{black!55}\rule[0.5ex]{0.25\linewidth}{0.4pt}\quad
  \textsc{Lecture #1}\quad\rule[0.5ex]{0.25\linewidth}{0.4pt}\hfill}\par\medskip}

% ---------- Notation ----------
\newcommand{\Obs}{\mathcal{O}}              % inertial observer / frame O
\newcommand{\Obar}{\bar{\mathcal{O}}}       % frame \bar O
\newcommand{\Ev}[1]{\mathcal{#1}}           % events: \Ev{A}, \Ev{B}, ...
\newcommand{\Lam}[2]{\Lambda^{#1}{}_{#2}}   % \Lam{i'}{j}
\newcommand{\dd}{\mathrm{d}}
\newcommand{\pd}[2]{\frac{\partial #1}{\partial #2}}
\newcommand{\dv}[2]{\frac{\dd #1}{\dd #2}}

% ---------- TikZ styles for spacetime diagrams ----------
\tikzset{
  axis/.style={-{Stealth[length=2mm]}, thick},
  paxis/.style={-{Stealth[length=2mm]}, thick, blue!60!black},   % axes of the moving frame
  light/.style={dashed, orange!80!black},
  wline/.style={very thick, red!70!black},
  event/.style={circle, fill=black, inner sep=1.2pt},
}
